\chapter{Partial Thromboplastin Time (PTT/aPTT)}

\section*{What Is This Test?}
The activated partial thromboplastin time (aPTT) measures the time required for plasma to clot after activation of the intrinsic and common coagulation pathways. It evaluates factors XII, XI, IX, VIII (intrinsic) and factors X, V, II, and I (common). The term ``partial'' refers to the use of phospholipid (partial thromboplastin) rather than complete tissue factor. The ``activated'' designation refers to the addition of a contact activator (silica, kaolin, ellagic acid) that standardizes activation of factor XII. The aPTT is used to screen for intrinsic pathway factor deficiencies (hemophilia A and B), to monitor unfractionated heparin therapy, and to detect lupus anticoagulant. It is one of the most commonly ordered coagulation tests, along with the PT/INR \cite{tierney2023current}.

\section*{Alternative Names}
aPTT; PTT; Activated Partial Thromboplastin Time; Partial Thromboplastin Time; APTT; Intrinsic Pathway Coagulation Test

\section*{Principle}
The aPTT is performed by adding a contact activator (silica, kaolin, ellagic acid, or celite) and phospholipid reagent (derived from rabbit brain, soybean, or synthetic sources) to citrated platelet-poor plasma. After incubation at 37\textdegree{}C for 3--5 minutes for optimal contact activation, calcium chloride is added. The contact activator initiates the intrinsic pathway: factor XII is activated to XIIa, which activates XI to XIa, which activates IX to IXa. Factor IXa, with factor VIIIa, phospholipid, and calcium (the tenase complex), activates factor X to Xa. Factor Xa, with factor Va, phospholipid, and calcium (the prothrombinase complex), converts prothrombin to thrombin, which cleaves fibrinogen to fibrin, forming a clot. Clot formation is detected optically or mechanically \cite{tietz2012textbook}. The aPTT is highly sensitive to heparin (heparin potentiates antithrombin, which inactivates thrombin and factor Xa). Therapeutic range for unfractionated heparin is typically 1.5--2.5 times control (60--100 seconds). Each laboratory establishes its therapeutic range by correlating aPTT with heparin anti-Xa activity (target 0.3--0.7 IU/mL) \cite{kaushansky2021williams}.

\section*{History}
The PTT was developed in the 1950s as a test for the intrinsic coagulation pathway. In 1961, Proctor and Rapaport introduced the activated PTT by adding kaolin as a contact activator, standardizing the contact phase and shortening clotting times. The aPTT became widely adopted for screening intrinsic pathway disorders, monitoring heparin therapy (replacing the less reliable Lee-White whole blood clotting time), and detecting lupus anticoagulant. Modern aPTT reagents use standardized phospholipid compositions and automated detection methods.

\section*{Why This Test Is Ordered}
\begin{itemize}
  \item Monitoring unfractionated heparin therapy --- aPTT is the standard test; therapeutic range 1.5--2.5 times control, corresponding to heparin anti-Xa of 0.3--0.7 IU/mL \cite{ponikowski2016esc}
  \item Screening for hemophilia A (factor VIII deficiency) and hemophilia B (factor IX deficiency) --- prolonged aPTT with normal PT
  \item Preoperative coagulation evaluation in patients with history of bleeding, easy bruising, or menorrhagia
  \item Evaluation of unexplained bleeding --- isolated prolonged aPTT suggests intrinsic pathway factor deficiency or von Willebrand disease
  \item Detection of lupus anticoagulant --- prolonged aPTT that fails to correct with mixing studies
  \item Diagnosis of DIC --- prolonged aPTT and PT with thrombocytopenia, low fibrinogen, elevated D-dimer
  \item Monitoring heparin during cardiac surgery, ECMO, or hemodialysis
  \item Screening before invasive procedures (lumbar puncture, epidural catheter, surgery)
\end{itemize}

\section*{Is This a Primary or Secondary Test}
The aPTT is a primary coagulation screening test for the intrinsic and common pathways, heparin monitoring, and lupus anticoagulant detection.

\section*{How This Test Is Performed}
A venous blood sample is collected in a light blue-top tube containing 3.2\% sodium citrate (9:1 blood-to-citrate ratio). The tube must be filled exactly to the marked line. The sample is centrifuged at 1,500--2,500 $\times$ g for 15 minutes to obtain platelet-poor plasma. On the automated analyzer, plasma is incubated with aPTT reagent at 37\textdegree{}C for 3--5 minutes, calcium chloride is added, and time to clot formation is measured optically or mechanically \cite{fischbach2023manual}. Venipuncture takes less than 5 minutes. Automated analysis is completed in 2--5 minutes.

\section*{What Preparation Is Needed}
No special preparation is required. Fasting is not necessary. The patient must disclose all medications, especially heparin (unfractionated or LMWH), warfarin, direct oral anticoagulants (rivaroxaban, apixaban prolong aPTT; dabigatran prolongs aPTT), and recent factor concentrate administration. Antiplatelet agents do not affect aPTT.

\section*{Contraindications and Precautions}
No absolute contraindications. Critical pre-analytical requirements: the light blue-top tube must be filled exactly to the line (9:1 ratio) --- underfilling causes falsely prolonged aPTT, overfilling risks clotting and consumption; samples must be gently inverted 3--4 times; hemolyzed, clotted, or severely lipemic specimens are unacceptable; traumatic venipuncture shortens aPTT; heparin contaminated samples produce markedly prolonged aPTT; polycythemia (hematocrit >55\%) reduces plasma volume, requiring an adjusted citrate tube; platelet-poor plasma is essential (residual platelets neutralize lupus anticoagulant); testing within 4 hours at room temperature; each laboratory must establish its own heparin therapeutic range; aPTT is not useful for monitoring low-molecular-weight heparin (LMWH) or DOACs.

\section*{Risks}
Minimal: mild pain, bruising, hematoma, lightheadedness, vasovagal syncope. Patients on anticoagulants have increased hematoma risk --- firm pressure for 3--5 minutes.

\section*{What Happens After the Test}
Results are available within 30--60 minutes. For heparin monitoring, aPTT is checked 6 hours after dose change. For prolonged aPTT evaluation: isolated prolonged aPTT (normal PT) suggests intrinsic pathway deficiency (factors VIII, IX, XI, XII), lupus anticoagulant, or heparin effect; combined prolonged aPTT and PT suggests common pathway deficiency (X, V, II, fibrinogen), vitamin K deficiency, warfarin, liver disease, or DIC. Mixing study distinguishes factor deficiency from inhibitor \cite{wallach2022interpretation}.

\section*{Understanding Results}

\subsection*{Below Normal (Shortened aPTT)}
\begin{itemize}
  \item \textbf{Elevated factor VIII (acute-phase response):} Acute illness, inflammation, pregnancy, surgery, trauma, malignancy. Associated with increased thrombotic risk.
  \item \textbf{Pregnancy and estrogen therapy:} Elevated factor VIII, factor VII, and vWF shorten aPTT.
  \item \textbf{Pre-analytical artifact:} Traumatic venipuncture, partial clotting, inadequate citrate (overfilled tube), polycythemia.
\end{itemize}

\subsection*{Above Normal (Prolonged aPTT)}
\begin{itemize}
  \item \textbf{Unfractionated heparin therapy:} Intentional prolongation to 1.5--2.5 times control.
  \item \textbf{Hemophilia A (factor VIII deficiency):} X-linked recessive. Factor VIII <1\% (severe --- spontaneous hemarthroses, deep muscle bleeding), 1--5\% (moderate), 5--40\% (mild --- post-traumatic/surgical bleeding). Markedly prolonged aPTT, normal PT, normal platelet count. Diagnosed by factor VIII assay.
  \item \textbf{Hemophilia B (factor IX deficiency):} X-linked recessive, identical presentation to hemophilia A. Diagnosed by factor IX assay.
  \item \textbf{Von Willebrand disease:} Most common inherited bleeding disorder (prevalence 1\%). vWF mediates platelet adhesion and stabilizes factor VIII. Decreased factor VIII prolongs aPTT. Mucocutaneous bleeding (epistaxis, gingival bleeding, menorrhagia, easy bruising, prolonged bleeding after dental work). Diagnosed by vWF antigen, vWF ristocetin cofactor activity, factor VIII activity, and vWF multimer analysis.
  \item \textbf{Hemophilia C (factor XI deficiency):} Autosomal recessive (Ashkenazi Jewish population). Milder bleeding, typically post-traumatic/post-surgical. Prolonged aPTT, normal PT.
  \item \textbf{Factor XII, prekallikrein, or HMWK deficiency:} Prolonged aPTT but NO clinical bleeding (paradoxically, may have thrombotic risk). Important to recognize to avoid unnecessary investigation.
  \item \textbf{Lupus anticoagulant (antiphospholipid antibody):} Acquired IgG/IgM autoantibody binding phospholipid-protein complexes. Prolongs phospholipid-dependent clotting tests (aPTT, dRVVT) in vitro but is associated with thrombosis (arterial and venous) and recurrent pregnancy loss in vivo, NOT bleeding. Diagnosed by: prolonged aPTT, failure to correct on 1:1 mixing with normal plasma (inhibitor pattern), correction or shortening on addition of excess phospholipid (hexagonal phase activation, platelet neutralization procedure), and positive dRVVT (dilute Russell viper venom time) \cite{devreese2018laboratory}.
  \item \textbf{Vitamin K deficiency, warfarin therapy:} Affect factors II, VII, IX, X --- prolonged aPTT and PT \cite{keeling2011guidelines}.
  \item \textbf{Liver disease:} Impaired synthesis of all coagulation factors except factor VIII; prolonged aPTT and PT.
  \item \textbf{DIC:} Consumptive coagulopathy with prolonged aPTT and PT, thrombocytopenia, low fibrinogen, elevated D-dimer, schistocytes.
  \item \textbf{Acquired factor inhibitors:} Acquired hemophilia A (autoantibody against factor VIII) presents with severe bleeding in adults without prior bleeding history, prolonged aPTT, factor VIII inhibitor by Bethesda assay.
  \item \textbf{Direct oral anticoagulants:} Dabigatran (direct thrombin inhibitor) prolongs aPTT; rivaroxaban and apixaban may slightly prolong aPTT.
  \item \textbf{Heparin contamination:} Collection from heparinized line.
\end{itemize}

\section*{Normal Results}
\begin{itemize}
  \item \textbf{Adults:} 25--35 seconds (reagent-dependent; each laboratory establishes its own reference range)
  \item \textbf{Therapeutic heparin range:} 1.5--2.5 times the mean of the normal reference range (typically 60--100 seconds), calibrated to heparin anti-Xa of 0.3--0.7 IU/mL
  \item \textbf{Newborns:} Prolonged aPTT at birth (40--60 seconds) due to low vitamin K-dependent factors and factor XII/prekallikrein; declines to adult levels by 3--6 months
  \item \textbf{Pregnancy:} Mildly shortened aPTT (third trimester) due to elevated factor VIII
  \item \textbf{Note:} aPTT reference ranges are highly reagent and instrument-specific; results cannot be directly compared between laboratories
\end{itemize}

\section*{Follow-up Tests}
\begin{itemize}
  \item \textbf{PT/INR:} To determine if prolongation is isolated to intrinsic pathway or involves both pathways
  \item \textbf{aPTT mixing study:} Patient plasma mixed 1:1 with normal pooled plasma; correction indicates factor deficiency; failure to correct after 2-hour incubation at 37\textdegree{}C indicates inhibitor (lupus anticoagulant or specific factor inhibitor)
  \item \textbf{Specific coagulation factor assays:} Factor VIII (hemophilia A, VWD), factor IX (hemophilia B), factor XI, factor XII, and others based on mixing study pattern
  \item \textbf{Factor VIII inhibitor (Bethesda assay):} When acquired hemophilia A is suspected
  \item \textbf{von Willebrand factor panel:} vWF antigen, vWF ristocetin cofactor activity, factor VIII activity, vWF multimers
  \item \textbf{Lupus anticoagulant panel:} dRVVT screen and confirm, hexagonal phase phospholipid neutralization (StaClot LA), silica clotting time
  \item \textbf{Anticardiolipin antibodies (IgG, IgM) and anti-beta-2-glycoprotein I antibodies:} For antiphospholipid syndrome diagnosis
  \item \textbf{Heparin anti-Xa assay:} For accurate heparin monitoring when aPTT is unreliable (lupus anticoagulant, baseline prolonged aPTT, high factor VIII)
  \item \textbf{Thrombin time (TT):} To detect heparin contamination, fibrinogen deficiency, dysfibrinogenemia, or fibrin degradation products
  \item \textbf{Fibrinogen (Clauss method):} When common pathway deficiency or DIC is suspected
  \item \textbf{D-dimer:} For DIC evaluation
  \item \textbf{Complete blood count with platelet count:} Assess for thrombocytopenia
  \item \textbf{Liver function tests:} When liver disease is suspected as cause of factor deficiencies
\end{itemize}

\section*{References}
\bibliographystyle{unsrtnat}
\bibliography{references}

\clearpage
